% The Collected Mathematical Papers of Arthur Cayley, Vol. XI
% Transcription of PDF pages 176--200 (printed pages 152--176)
% Papers:
%   745 (continued) ON THE SCHWARZIAN DERIVATIVE, AND THE POLYHEDRAL FUNCTIONS
%                                                                    (PDF pp176--200, printed 152--176)
% Scan source: collectedmathema11cayluoft.pdf

\documentclass[11pt]{article}
\usepackage[a4paper,margin=0.65in]{geometry}
\usepackage[T1]{fontenc}
\usepackage{amsmath,amssymb,amsfonts}
\usepackage{graphicx}
\usepackage{pdfpages}
\usepackage[english,shorthands=off]{babel}
\usepackage{fancyhdr}
\usepackage[bookmarks=false,colorlinks=false]{hyperref}
\usepackage[protrusion=true,expansion=false]{microtype}
\emergencystretch=3em
\tolerance=600
\providecommand{\modd}{\mathrm{mod}}
\providecommand{\Disc}{\mathrm{Disc}}
\providecommand{\canont}{}
\providecommand{\asterism}{*}
\providecommand{\Hessian}{H}
\providecommand{\tome}{}
\providecommand{\page}{p.}
\pagestyle{fancy}\fancyhf{}\fancyhead[L]{Pages 176--200}\fancyhead[R]{Cayley --- Collected Papers, Vol. XI}\renewcommand{\headrulewidth}{0.4pt}
\setlength{\headheight}{15pt}

\newcommand{\artno}[1]{\textbf{#1.}\quad}

\begin{document}

% ============================================================
% PDF p176 (printed 152) --- Paper 745 continued
% ============================================================

Again differentiating, we have
\[
   \frac{s'''}{s'} - \left(\frac{s''}{s'}\right)^{\!2}
   = S'^{2}\!\left[\frac{s_{3}}{s_{1}} - \left(\frac{s_{2}}{s_{1}}\right)^{\!2}\right]
   + 3S'\frac{s_{2}}{s_{1}}\,\frac{S''}{S'} - \left(\frac{S''}{S'}\right)^{\!2}.
\]
But
\[
   -\tfrac{1}{2}\!\left(\frac{s''}{s'}\right)^{\!2} = S'^{2}\!\left[-\tfrac{1}{2}\!\left(\frac{s_{2}}{s_{1}}\right)^{\!2}\right] - S'\frac{s_{2}}{s_{1}}\,\frac{S''}{S'} - \tfrac{1}{2}\!\left(\frac{S''}{S'}\right)^{\!2},
\]
and consequently
\[
   \frac{s'''}{s'} - \tfrac{3}{2}\!\left(\frac{s''}{s'}\right)^{\!2}
   = S'^{2}\!\left[\frac{s_{3}}{s_{1}} - \tfrac{3}{2}\!\left(\frac{s_{2}}{s_{1}}\right)^{\!2}\right]
   + \frac{S'''}{S'} - \tfrac{3}{2}\!\left(\frac{S''}{S'}\right)^{\!2};
\]
that is,
\[
   \{s,\,x\} = \left(\frac{dS}{dx}\right)^{\!2}\{s,\,S\} + \{S,\,x\},
\]
the required formula.

In a very similar manner, taking $x$ a function of $X$, it is shown that
\[
   \{s,\,x\} = \left(\frac{dX}{dx}\right)^{\!2}\{s,\,X\} - \{x,\,X\}.
\]

\artno{3} If in this formula we write $S$ for $s$, and substitute the resulting value of
$\{S,\,x\}$ in the former formula, we have
\[
   \{s,\,x\} = \left(\frac{dS}{dx}\right)^{\!2}\{s,\,S\} + \left(\frac{dX}{dx}\right)^{\!2}\{S,\,X\} - \{x,\,X\},
\]
which is the formula for the change of both variables. It in fact includes the other
two: viz.\ writing $X = x$, or $S = x$, and observing that $\{x,\,x\} = 0$, we have the
other two formul\ae.

\artno{4} By putting in the first formula $S = x$, we obtain
\[
   \{x,\,s\} = -\left(\frac{ds}{dx}\right)^{\!2}\{s,\,x\},
\]
a formula for the interchange of the variables.

\artno{5} Writing $S = \dfrac{as+b}{cs+d}$, and using for a moment the accents to denote differentiation
in regard to $s$, we have
\[
   S' = \frac{ad-bc}{(cs+d)^{2}}, \qquad \frac{S''}{S'} = \frac{-2c}{cs+d},
\]
and thence
\[
   \frac{S'''}{S'} - \left(\frac{S''}{S'}\right)^{\!2} = \frac{2c^{2}}{(cs+d)^{2}},
\]
\[
   -\tfrac{1}{2}\!\left(\frac{S''}{S'}\right)^{\!2} = \frac{-2c^{2}}{(cs+d)^{2}}.
\]
Consequently $\{S,\,s\} = 0$, whence also $\{s,\,S\} = 0$.

\newpage
% ============================================================
% PDF p177 (printed 153)
% ============================================================
Hence in the first formula $\{S,\,x\} = \{s,\,x\}$, that is,
\[
   \left\{\frac{as+b}{cs+d},\,x\right\} = \{s,\,x\};
\]
viz.\ we may, in the derivative $\{s,\,x\}$, write for $s$ any homographic function $\dfrac{as+b}{cs+d}$
of $s$.

\artno{6} Again, if $X = \dfrac{\alpha x+\beta}{\gamma x+\delta}$, then from the second formula
\[
   \{s,\,x\} = \frac{(\alpha\delta - \beta\gamma)^{2}}{(\gamma x+\delta)^{4}}\{s,\,X\};
\]
that is,
\[
   \left\{s,\,\frac{\alpha x+\beta}{\gamma x+\delta}\right\} = \frac{(\gamma x+\delta)^{4}}{(\alpha\delta - \beta\gamma)^{2}}\{s,\,x\};
\]
and here, changing $s$ into $(as+b) \div (cs+d)$, we have finally
\[
   \left\{\frac{as+b}{cs+d},\,\frac{\alpha x+\beta}{\gamma x+\delta}\right\} = \frac{(\gamma x+\delta)^{4}}{(\alpha\delta - \beta\gamma)^{2}}\{s,\,x\},
\]
which is the formula for the homographic transformation of the two variables $s$, $x$.

\artno{7} Let $s$ be a given function of $x$, the equation $\{S,\,x\} = \{s,\,x\}$ is a differential
equation of the third order in $S$, and by what precedes, its general integral is $S = \dfrac{as+b}{cs+d}$.

The direct process is as follows: we have a first integral $\dfrac{S''}{S'} = \dfrac{s''}{s'} - \dfrac{2cs'}{cs+d}$; a second
integral $\log S' = \log s' - 2\log(cs+d) + \mathrm{const.}$, that is, $S' = \dfrac{cAs'}{(cs+d)^{2}}$; and thence a final
integral $S = B - \dfrac{A}{cs+d}$, which is equivalent to the foregoing value of $S$.

\bigskip

\begin{center}
\textit{The Quadric Function of three or more Inverts. Art.\ Nos.\ 8 to 15.}
\end{center}

\medskip

\artno{8} We consider a quadric function of any number of inverts $\dfrac{1}{x-a},\, \dfrac{1}{x-\beta},\, \ldots,$ all
of them different: it is assumed that the constant term is $=0$, and also that the
sum of the coefficients of the linear terms is $=0$. We have therefore square terms
$\dfrac{a}{(x-a)^{2}}, \ldots,$ product terms $\dfrac{h}{x-a \,.\, x-\beta},$ and linear terms $\dfrac{A}{x-a},$ where the
sum of the coefficients $A$ is $=0$. Any product term $\dfrac{h}{x-a \,.\, x-\beta}$ is expressible in the form of a
difference $\dfrac{h}{a-\beta}\!\left(\dfrac{1}{x-a} - \dfrac{1}{x-\beta}\right)$ of two linear terms, and (the coefficients of these

\newpage
% ============================================================
% PDF p178 (printed 154)
% ============================================================
being equal), after it is thus expressed, the sum of the coefficients of the linear terms
is still $=0$. The function is thus always expressible in the form
\[
   \frac{a}{(x-a)^{2}} + \frac{b}{(x-\beta)^{2}} + \cdots + \frac{A}{x-a} + \frac{B}{x-\beta} + \cdots,
\]
where the sum $A + B + \cdots$ is $=0$: this may be called the reduced form.

\artno{9} Observe that any particular invert $\dfrac{1}{x-a}$ may disappear altogether from the
reduced form: this will be the case if $a = 0$, that is, if the original form contains no
term in $\dfrac{1}{(x-a)^{2}}$, and if also $A = 0$. An invert thus disappearing from the reduced
form is said to be non-essential: and the inverts which do not disappear are said to
be essential. The original form contains in appearance the non-essential inverts, but
it is really a quadric function of the essential inverts only.

\artno{10} Imagine the original function expressed as a rational fraction, the denominator
being the product $(x-a)^{2}(x-\beta)^{2}(x-\gamma)^{2}\cdots$ of the squared factors corresponding to all
the inverts (non-essential as well as essential): the numerator will be in general of a
degree less by 2 than that of the denominator, but the coefficients of any one or
more of the higher powers of $x$ may vanish, and the numerator will then be of a
lower degree. But this numerator will for any non-essential invert $\dfrac{1}{x-\gamma}$ contain the
factor $(x-\gamma)^{2}$, or, dividing the numerator and denominator each by this factor, the
difference of the degrees of the numerator and denominator will remain unaltered;
that is, the difference will have the same value whether we do or do not attend to
the non-essential inverts; or say it will have the same value for the original form and
for the reduced form.

\artno{11} It is to be remarked that the linear terms $\dfrac{A}{x-a} + \dfrac{B}{x-\beta} + \dfrac{C}{x-\gamma} + \cdots$, where
$A + B + C + \cdots = 0$, can be (and that in a variety of ways) expressed as a sum of
differences $\dfrac{1}{x-a} - \dfrac{1}{x-\beta}$, that is, as a sum of product-terms $\dfrac{1}{x-a \,.\, x-\beta}$. Hence the
quadric function can be (and that in a variety of ways) expressed as a homogeneous
function $(a, b, \ldots \,\mid\!\!\cdot\!\!\cdot\!\!\cdot\,\mid \frac{1}{x-a}, \frac{1}{x-\beta}, \ldots)^{\!2}$; we must have in the form all the essential inverts,
and we need have these only. Supposing that this is so, and that the number of
the essential inverts is $=n$, then the number of constants is $=\tfrac{1}{2}n(n+1)$, whereas the
number of constants in the reduced form is only $=2n-1$: hence the coefficients are
not determinate; or, what is the same thing, we may have different quadric functions
having each of them the same reduced function; these quadric functions, as having
the same reduced function, can only differ by multiples of the evanescent expressions
\[
   \frac{\beta-\gamma}{x-\beta \,.\, x-\gamma} + \frac{\gamma-a}{x-\gamma \,.\, x-a} + \frac{a-\beta}{x-a \,.\, x-\beta}, \quad \&\text{c.}
\]

\newpage
% ============================================================
% PDF p179 (printed 155)
% ============================================================
In particular, if the number of essential inverts is $=3$, then the quadric function is
of the form
\[
   \left(a,\,b,\,c,\,f,\,g,\,h \,\middle|\!\!\!\!\mid \frac{1}{x-a},\,\frac{1}{x-\beta},\,\frac{1}{x-\gamma}\right)^{\!2},
\]
which contains one superfluous constant, and equivalent functions differ only by a
multiple of
\[
   \frac{\beta-\gamma}{x-\beta \,.\, x-\gamma} + \frac{\gamma-a}{x-\gamma \,.\, x-a} + \frac{a-\beta}{x-a \,.\, x-\beta}.
\]

\artno{12} A quadric function such that the degree of the numerator is less by 4 than
that of the denominator is said to be ``curtate.''

The conditions, in order that the function
\[
   \left(a,\,b,\,c,\,f,\,g,\,h \,\middle|\!\!\!\!\mid \frac{1}{x-a},\,\frac{1}{x-\beta},\,\frac{1}{x-\gamma}\right)^{\!2}
\]
may be curtate, are easily found to be
\[
   a + b + c + 2f + 2g + 2h = 0,
\]
\[
   a(a+h+g) + \beta(h+b+f) + \gamma(g+f+c) = 0;
\]
and by reason of the superfluous constant we are at liberty to assume a third condition:
the three conditions may be taken to be $a+h+g$, $h+b+f$, $g+f+c$ each $=0$; and
this being so the values of $f$, $g$, $h$ are $=\tfrac{1}{2}(a-b-c)$, $\tfrac{1}{2}(-a+b-c)$, $\tfrac{1}{2}(-a-b+c)$
respectively. Hence the form is
\[
   \left(a,\,b,\,c,\,\tfrac{1}{2}(a-b-c),\,\tfrac{1}{2}(-a+b-c),\,\tfrac{1}{2}(-a-b+c) \,\middle|\!\!\!\!\mid \frac{1}{x-a},\,\frac{1}{x-\beta},\,\frac{1}{x-\gamma}\right)^{\!2},
\]
which, as already mentioned, we denote by
\[
   \left(a,\,b,\,c \,.{}^{\!\cdot\!\cdot\!\cdot}\,\middle|\!\!\!\!\mid \frac{1}{x-a},\,\frac{1}{x-\beta},\,\frac{1}{x-\gamma}\right)^{\!2}.
\]

We have thus the theorem that a curtate function of any number of inverts, but with
only the three essential inverts
\[
   \frac{1}{x-a},\,\frac{1}{x-\beta},\,\frac{1}{x-\gamma},
\]
is always expressible in the foregoing form
\[
   \left(a,\,b,\,c \,.{}^{\!\cdot\!\cdot\!\cdot}\,\middle|\!\!\!\!\mid \frac{1}{x-a},\,\frac{1}{x-\beta},\,\frac{1}{x-\gamma}\right)^{\!2}.
\]

\artno{13} It may be remarked that the function $(a,\,b,\,c \,.{}^{\!\cdot\!\cdot\!\cdot}\,\mid X,\,Y,\,Z)^{2}$ is a function of
the differences of the variables $X$, $Y$, $Z$; and similarly, in the case of four variables,
a function $(a,\,b,\,c,\,d,\,f,\,g,\,h,\,l,\,m,\,n \,\mid X,\,Y,\,Z,\,W)^{2}$, for which
\[
   a+h+g+l,\ \ h+b+f+m,\ \ g+f+c+n,\ \ l+m+n+d,
\]

\newpage
% ============================================================
% PDF p180 (printed 156)
% ============================================================
are each $=0$, is a function of the differences of the variables $X$, $Y$, $Z$, $W$: and so in
general. Any such function is said to be ``diaphoric'': and it is easy to see that,
taking for the variables any inverts whatever, a diaphoric function is always curtate.

\artno{14} The function
\[
   \left(\cdots a \cdots,\,\cdots b \cdots,\,\cdots c \cdots \,\middle|\!\!\!\!\mid \frac{1}{x-a},\,\frac{1}{x-a_{1}},\,\cdots\right)^{\!2},
\]
where the coefficients $a, b, c, \ldots$ satisfy the relation $a + b + c + \cdots = -2$, is diaphoric,
and therefore curtate. In fact, forming the sum, coeff.\ $\dfrac{1}{(x-a)^{2}} + \tfrac{1}{2}\,\text{coeff.}\,\dfrac{1}{x-a\,.\,x-\beta} + \cdots$,
this is $=a + \tfrac{1}{2}a^{2} - \tfrac{1}{2}ab - \tfrac{1}{2}ac - \cdots$, $=-\tfrac{1}{2}a(2 + a + b + c + \cdots)$, which is $=0$; and similarly
the other conditions are satisfied.

\artno{15} The function
\[
   \left(\cdots a \cdots,\,\cdots b \cdots,\,\cdots c \cdots \,\middle|\!\!\!\!\mid \frac{1}{x-a},\,\frac{1}{x-a_{1}},\,\cdots,\,\frac{1}{x-\beta},\,\cdots\right)^{\!2}
\]
regarded as a function of the inverts $\dfrac{1}{x-a},\,\dfrac{1}{x-a_{1}},\,\cdots,\,\dfrac{1}{x-\beta},\,\cdots,$ where
\[
   a + a_{1} + \cdots = b + b_{1} + \cdots = c + c_{1} + \cdots, = h \text{ suppose,}
\]
is diaphoric, and therefore curtate. In fact, the condition in regard to $\dfrac{1}{x-a}$ is
\[
   a(a + a_{1} + \cdots) + \tfrac{1}{2}(a - b + c)(a c + a c_{1} + \cdots) + \tfrac{1}{2}(-a + b - c)(a b + a b_{1} + \cdots) + \tfrac{1}{2}(-a - b + c)(\cdots) = 0;
\]
that is,
\[
   ah\,\{a + \tfrac{1}{2}(-a + b - c) + \tfrac{1}{2}(-a - b + c)\} = 0,
\]
which is satisfied. And similarly the other conditions are satisfied.

\bigskip

\begin{center}
\textit{The functions $P$, $Q$, $R$. Art.\ Nos.\ 16 to 20.}
\end{center}

\medskip

\artno{16} We consider $P$, $Q$, $R$, rational and integral functions of $z$, such that $P + Q + R = 0$:
hence, using the accent to denote differentiation in regard to $z$, we have also $P' + Q' + R' = 0$;
and therefore $QR' - Q'R = RP' - R'P = PQ' - P'Q$, $= \Theta$ suppose: and we require to find
$P$, $Q$, $R$, such that the function $\Theta$ contains only the factors of $P$, $Q$, $R$.

\artno{17} It is to be observed that, effecting upon a solution $P$, $Q$, $R$ any linear substitution
$\div (\alpha z + \beta) \div (\gamma z + \delta)$, and omitting the common denominator, we have a solution;
but this is regarded as identical with the original solution. The three functions, if

\newpage
% ============================================================
% PDF p181 (printed 157)
% ============================================================
not originally of the same order, can thus be made to be of the same order; or by
taking account of the root $z = \infty$, we may in the original case regard them as being
of the same order, and it is convenient so to regard them: say they are taken to
be of the same order $\delta$. And there is clearly no loss of generality in taking the
three functions to be prime to each other; for any common factor of two of them
would divide the third, and might therefore be struck out.

\artno{18} We may therefore write
\[
   P = F\Pi(z-l)^{p}, \quad Q = G\Pi(z-m)^{q}, \quad R = H\Pi(z-n)^{r},
\]
where $(z-l)^{p}$ is taken to denote the distinct simple or multiple factors of $P$, and
the like as regards $Q$ and $R$; the factors $z-l,\,z-m,\,z-n$ are thus all of them different.
And we have $\delta = \Sigma p$, $= \Sigma q$, $= \Sigma r$.

\artno{19} It is at once seen that $\Theta$ is of the degree $2\delta - 2$, and moreover that it
contains the factors $\Pi(z-l)^{p-1}$, $\Pi(z-m)^{q-1}$, $\Pi(z-n)^{r-1}$; hence it contains the factor
\[
   \Pi(z-l)^{p-1}(z-m)^{q-1}(z-n)^{r-1}.
\]
Suppose the number of distinct indices $p$ is $=\sigma_{1}$, that of distinct indices $q$ is $\sigma_{2}$, and
that of distinct indices $r$ is $\sigma_{3}$; then the degree of the factor is $=3\delta - \sigma_{1} - \sigma_{2} - \sigma_{3}$;
and if this be $=2\delta - 2$, then $\Theta$ can have no other valuable factor: viz.\ if the numbers
$\sigma_{1}, \sigma_{2}, \sigma_{3}$ of the distinct indices $p, q, r$ respectively are such that $\sigma_{1} + \sigma_{2} + \sigma_{3} = \delta + 2$,
a relation which is henceforth taken to be satisfied, then we have
\[
   \Theta = K\Pi(z-l)^{p-1}(z-m)^{q-1}(z-n)^{r-1}.
\]
As already in effect remarked, the conclusion extends to the case where $P$, $Q$, $R$ are
not of the same degree; the equation $P + Q + R = 0$ here implies that two functions,
say $P$, $Q$, are of the same degree, and the third function $R$ of an inferior degree;
but, this being so, we have only to regard $R$ as containing the factor $\left(1 - \dfrac{z}{\infty}\right)^{\!t}$ of
the degree $t$ proper for raising its degree up to that of $P$ or $Q$.

\artno{20} Solutions are given in the following PQR-Table; in which, where required,
the proper factor $\left(1 - \dfrac{z}{\infty}\right)$ has been added; the first column headed Ref.\ No.\ (Reference
Number) will be explained further on. The Annex to the same Table will also be
explained.

\newpage
% ============================================================
% PDF p182 (printed 158) --- THE PQR-TABLE (rotated landscape in original)
% ============================================================
\begin{center}
\textsc{The PQR-Table.}
\end{center}
\medskip
\noindent\resizebox{\textwidth}{!}{%
\begin{tabular}{|c|l|l|l|l|}
\hline
Ref. No. & $P =$ & $Q =$ & $R =$ & $\Theta =$ \\
\hline
1   & $z^{n}$
    & $-1\!\left(1 - \dfrac{z}{\infty}\right)^{\!n}$
    & $-z^{n} + 1$
    & $n z^{n-1}\!\left(1 - \dfrac{z}{\infty}\right)^{\!n-1}$ \\[6pt]
2   & $4 z^{n}\!\left(1 - \dfrac{z}{\infty}\right)^{\!n}$
    & $-(z^{n}+1)^{2}$
    & $(z^{n}-1)^{2}$
    & $-4 n z^{n-1}(z^{n}+2)(z^{n}-1)\!\left(1 - \dfrac{z}{\infty}\right)^{\!n-1}$ \\[6pt]
$3\tfrac{1}{2}$ & $(z^{4} + 2\sqrt{-3}\,z^{2} + 1)^{3}$
    & $-12\sqrt{-3}\,z^{2}(z^{4}-1)^{2}\!\left(1 - \dfrac{z}{\infty}\right)^{\!2}$
    & $-(z^{4} - 2\sqrt{-3}\,z^{2} + 1)^{3}$
    & $24\sqrt{-3}\,(z^{4} + 2\sqrt{-3}\,z^{2} + 1)^{2}$ \\
    & & & & $(z^{4} - 2\sqrt{-3}\,z^{2} + 1)^{2}\,z(z^{4}-1)\!\left(1 - \dfrac{z}{\infty}\right)$ \\[6pt]
4   & $(z^{8} + 14 z^{4} + 1)^{3}$
    & $-(z^{12} - 33 z^{8} - 33 z^{4} + 1)^{2}$
    & $-108\,(z^{5} - z)^{4}\!\left(1 - \dfrac{z}{\infty}\right)^{\!4}$
    & $216\,(z^{8} + \cdots)^{2}(z^{12} - \cdots)\,(z^{5}-z)^{3}\!\left(1 - \dfrac{z}{\infty}\right)^{\!3}$ \\[6pt]
5   & $(z^{20} - 228 z^{15} + 494 z^{10} + 228 z^{5} + 1)^{3}$
    & $-(z^{30} - 522 z^{25} - 10005 z^{20} + 0 z^{15} - 10005 z^{10} + 522 z^{5} - 1)^{2}$
    & $-1728\,(z^{11} + 11 z^{6} - z)^{5}\!\left(1 - \dfrac{z}{\infty}\right)^{\!5}$
    & $-8640\,(z^{20} - \cdots)^{2}(z^{11} + \cdots)^{4}\!\left(1 - \dfrac{z}{\infty}\right)^{\!4}$ \\[6pt]
\hline
III, V, VII, VIII & $4 z\!\left(1 - \dfrac{z}{\infty}\right)$
    & $-(z+1)^{2}$
    & $(z-1)^{2}$
    & $-4(z+1)(z-1)$ \\[6pt]
IX  & $(z-4)^{3}$
    & $-(z-1)(z+8)^{2}$
    & $27\,z^{2}\!\left(1 - \dfrac{z}{\infty}\right)$
    & $27\,(z-4)^{2}(z+8)\,z$ \\[6pt]
X   & $z(z+8)^{3}$
    & $-(z^{2} - 20z - 8)^{2}$
    & $-64\,(z-1)^{3}\!\left(1 - \dfrac{z}{\infty}\right)$
    & $-64\,(z+8)^{2}(z^{2} - 20z\cdots)(z-1)$ \\[6pt]
XI  & $4(z^{2} - z + 1)^{3}$
    & $-(2 z^{3} - 3 z^{2} - 3 z + 2)^{2}$
    & $-27\,z^{2}(z-1)^{2}\!\left(1 - \dfrac{z}{\infty}\right)^{\!2}$
    & $-108\,(z^{2} - \cdots)^{2}(2 z^{3} - 3 z^{2} - \cdots)$ \\
    & & & & $\cdot\, z(z-1)\!\left(1 - \dfrac{z}{\infty}\right)$ \\[6pt]
XII & $z^{3}(z+5)^{2}(z+8)$
    & $-(z^{3} + 9 z^{2} + 12 z - 8)^{2}$
    & $-64\,(3z-1)\!\left(1 - \dfrac{z}{\infty}\right)^{\!5}$
    & $-960\,z^{2}(z+5)(z^{3} + 9 z^{2}\cdots)\!\left(1 - \dfrac{z}{\infty}\right)^{\!4}$ \\[6pt]
XIII & $(z^{2} + 14 z + 1)^{3}$
    & $-(z^{3} - 33 z + 1)^{2}$
    & $-108\,z(z-1)^{4}$
    & $-108\,(z^{2} + 14 z \cdots)^{2}(z^{3} - 33 z \cdots)(z-1)^{3}$ \\[6pt]
XIV & $(64 z + 189)(64 z^{2} + 133 z + 49)^{3}$
    & $-z\,(4096 z^{3} + 18816 z^{2} + 25725 z + \cdots)^{2}$
    & $-27 \cdot 7^{7}\,(z+1)^{3}\!\left(1 - \dfrac{z}{\infty}\right)^{\!5}$
    & $-135 \cdot 7^{7}\,(64 z^{2} + \cdots)^{2}(4096 z^{3} + \cdots)$ \\
    & & & & $\cdot\,(z+1)\!\left(1 - \dfrac{z}{\infty}\right)^{\!4}$ \\[6pt]
XV  & $-(5z - 27)(125 z^{3} - 25 z^{2} - 265 z - 243)^{3}$
    & $-(-3125 z^{5} + 9375 z^{4} + 18750 z^{3} + 8750 z^{2} + 30750 z + \cdots)^{2}$
    & $+1\,382\,400\,000\,z^{3}(z+1)^{2}\!\left(1 - \dfrac{z}{\infty}\right)^{\!5}$
    & $-1\,382\,400\,000\,(125 z^{3} - \cdots)^{2}\!\left(-3125 z^{5} + \cdots\right)\!z^{2}(z+1)\!\left(1 - \dfrac{z}{\infty}\right)^{\!4}$ \\
\hline
\end{tabular}%
}

\medskip
\noindent\small In the second half of the table the functions $P$, $Q$, $R$, except in the lines XII, XIV
and XV, were calculated by Brioschi; those for the lines XII and XIV were calculated
by Klein, but as regards line XV there would seem to have been some error at the
beginning of the calculation, and the values found by him are erroneous.\normalsize

\newpage
% ============================================================
% PDF p183 (printed 159) --- ANNEX TO THE PQR-TABLE
% ============================================================
\begin{center}
\textsc{Annex to the PQR-Table.}
\end{center}
\medskip
\noindent\resizebox{\textwidth}{!}{%
\begin{tabular}{|c|l|l|l|l|c|}
\hline
Ref. No. & $(a, b, c)^{*}$ & \multicolumn{3}{c|}{Calculation of $(ap)$, $(bq)$, $(cr)$.} & Polygon \\
\hline
1  & $\tfrac{1}{2}(1 - n^{-2}),\ \tfrac{1}{2}(1 - n^{-2}),\ 0$    & $an = 0$ & $bn = 0$ & $c1 = 0$ & I \quad Polygon \\
2  & $\tfrac{1}{2}(1 - n^{-2}),\ \tfrac{3}{8},\ \tfrac{3}{8}$     & $an = 0$ & $b2 = 0$ & $c2 = 0$ & I \quad Double Pyramid \\
3  & $\tfrac{4}{9},\ \tfrac{3}{8},\ \tfrac{4}{9}$                 & $a3 = 0$ & $b2 = 0$ & $c3 = 0$ & II \quad Tetrahedron \\
4  & $\tfrac{4}{9},\ \tfrac{3}{8},\ \tfrac{15}{32}$               & $a3 = 0$ & $b2 = 0$ & $c4 = 0$ & IV \quad Cube and Octahedron \\
5  & $\tfrac{4}{9},\ \tfrac{3}{8},\ \tfrac{12}{25}$               & $a3 = 0$ & $b2 = 0$ & $c5 = 0$ & VI \quad Dodecahedron and Icosahedron \\
\hline
III   & $\tfrac{4}{9},\ \tfrac{3}{8},\ \tfrac{4}{9}$              & $a1 = \tfrac{4}{9},\ a1 = \tfrac{4}{9}$ & $b2 = 0$ & $c2 = \tfrac{5}{18}\!:\!\Sigma\!\left(\tfrac{1}{z},\,\tfrac{1}{z-\infty},\,\tfrac{1}{z-1}\right)^{2}$ & Tetrahedron \\
V     & $\tfrac{15}{32},\ \tfrac{3}{8},\ \tfrac{4}{9}$            & $a1 = \tfrac{15}{32},\ a1 = \tfrac{15}{32}$ & $b2 = 0$ & $c2 = \tfrac{5}{18}\!:\,\&\text{c.}$ & Cube and Octahedron \\
VII   & $\tfrac{4}{9},\ \tfrac{3}{8},\ \tfrac{12}{25}$            & $a1 = \tfrac{4}{9},\ a1 = \tfrac{4}{9}$ & $b2 = 0$ & $c2 = \tfrac{21}{50}$ & Dodecahedron and Icosahedron \\
VIII  & $\tfrac{12}{25},\ \tfrac{3}{8},\ \tfrac{4}{9}$            & $a1 = \tfrac{12}{25},\ a1 = \tfrac{12}{25}$ & $b2 = 0$ & $c2 = \tfrac{5}{18}$ & ,, \\
IX    & $\tfrac{4}{9},\ \tfrac{3}{8},\ \tfrac{12}{25}$            & $a3 = 0$ & $b1 = \tfrac{3}{8},\ b2 = 0$ & $c1 = \tfrac{12}{25},\ c2 = \tfrac{21}{50}$ & ,, \\
X     & ,,\ \ ,,\ \ ,,                                            & $a1 = \tfrac{4}{9},\ a3 = 0$ & $b2 = 0$ & $c1 = \tfrac{12}{25},\ c3 = \tfrac{21}{50}$ & ,, \\
XI    & ,,\ \ ,,\ \ ,,                                            & $a3 = 0$ & $b2 = 0$ & $c2 = \tfrac{21}{50},\ c2 = \tfrac{21}{50}$ & ,, \\
XII   & ,,\ \ ,,\ \ ,,                                            & $a1 = \tfrac{4}{9},\ a2 = \tfrac{5}{18},\ a3 = 0$ & $b2 = 0$ & $c1 = \tfrac{12}{25},\ c5 = 0$ & ,, \\
XIII  & ,,\ \ ,,\ \ ,,                                            & $a3 = 0$ & $b2 = 0$ & $c1 = \tfrac{12}{25},\ c1 = \tfrac{12}{25}$ & ,, \\
XIV   & ,,\ \ ,,\ \ ,,                                            & $a1 = \tfrac{4}{9},\ a3 = 0$ & $b1 = \tfrac{3}{8},\ b2 = 0$ & $c2 = \tfrac{21}{50},\ c5 = 0$ & ,, \\
XV    & ,,\ \ ,,\ \ ,,                                            & $a1 = \tfrac{4}{9},\ a3 = 0$ & $b2 = 0$ & $c2 = \tfrac{21}{50},\ c3 = \tfrac{21}{50}$ & ,, \\
\hline
\end{tabular}%
}

\medskip
\noindent\small * Observe, as regards the $(a, b, c)$ column, that the line III agrees with 3; line V with 4 (only there is a transposition of $\tfrac{4}{9}$ and $\tfrac{15}{32}$); lines VII and VIII agree each of them with 5 (only as regards VIII there is a transposition of $\tfrac{4}{9}$ and $\tfrac{12}{25}$): the remaining lines IX to XV agree each of them with 5. Observe also as regards the Annex generally, the Roman numbering of the lines 2, 3, 4, 5; the lines after the first have thus the Roman numbers I to XV, corresponding to the Roman numbers used by Schwarz.\normalsize

\newpage
% ============================================================
% PDF p184 (printed 160)
% ============================================================
\begin{center}
\textit{The Differential Equations $\{x,\,z\}$ and $\{s,\,x\}$. Art.\ Nos.\ 21 to 45.}
\end{center}

\medskip

\artno{21} In reference to what follows, it is convenient to put $P = XP_{0}$, $P' = X_{1}P_{1}$,
$= P''$ multiplied by the product $\Pi(z-l)$ of the several factors taken each with the index
unity; and so for $Q$ and $R$: viz.\ we write
\[
   P,\, Q,\, R = XP_{0},\, YQ_{0},\, ZR_{0},
\]
\[
   P',\, Q',\, R' = X_{1}P_{0},\, Y_{1}Q_{0},\, Z_{1}R_{0},
\]
and the foregoing value of $\Theta$ then is
\[
   \Theta = KP_{0}Q_{0}R_{0}.
\]

We come now to the investigation of the leading theorem. Take $a$, $b$, $c$ arbitrary,
$f$, $g$, $h = b-c,\, c-a,\, a-b$; $P$, $Q$, $R$ functions of $z$ as above; and write
\[
   f(x-a) : g(x-b) : h(x-c) = P : Q : R,
\]
equations, which are consistent with each other and determine $x$ as a rational function
of $z$. Using, as before, the accent to denote differentiation in regard to $z$, and taking
the coefficients $(a, b, c)$ arbitrary, it is required to find the value of
\[
   \{x,\,z\} + x'^{2}\!\left(a, b, c \,.{}^{\!\cdot\!\cdot\!\cdot}\,\middle|\!\!\!\mid \frac{1}{x-a},\,\frac{1}{x-b},\,\frac{1}{x-c}\right)^{\!2}.
\]

\artno{22} Calculation of the first term $\{x,\,z\}$.

We have $x$ a function $(\alpha P + \beta R) \div (\gamma P + \delta R)$, $=\dfrac{P}{R}$ suppose; and thence $\{x,\,z\} = \left\{\dfrac{P}{R},\,z\right\}$;
for a moment; then
\[
   \frac{x'}{x} = \frac{(P/R)'}{P/R} = \frac{PR' - P'R}{PR} = -\frac{P_{0}Q_{0}\cdot KR_{0}\cdot}{Z\,Z_{1}\,R_{0}\,\cdot\,P/R} = -\frac{P_{1}Q_{0}}{Z_{1}\,R_{0}}.
\]

Substituting the values
\[
   P_{0} = \Pi(z-l)^{p-1},\quad Q_{0} = \Pi(z-m)^{q-1},\quad R_{0} = \Pi(z-n)^{r-1},\quad Z = \Pi(z-n),
\]
we have
\[
   \frac{x'}{x} = \Sigma\,\frac{p-1}{z-l} + \Sigma\,\frac{q-1}{z-m} - \Sigma\,\frac{r+1}{z-n},
\]
and thence
\[
   \{x,\,z\} = \Sigma\,\frac{p-1}{(z-l)^{2}} + \Sigma\,\frac{q-1}{(z-m)^{2}} - \Sigma\,\frac{r+1}{(z-n)^{2}}
\]
\[
   - \tfrac{1}{2}\!\left[\Sigma\,\frac{p-1}{z-l} + \Sigma\,\frac{q-1}{z-m} - \Sigma\,\frac{r+1}{z-n}\right]^{2},
\]
or say
\[
   \{x,\,z\} = \tfrac{1}{2}\!\left(\Sigma\,\frac{p_{1}-1}{(z-l)^{2}} + \Sigma\,\frac{p_{1}-1}{(z-l_{1})^{2}} - \Sigma\,\frac{q-1}{(z-m)^{2}} - \Sigma\,\frac{q-1}{(z-m_{1})^{2}} + \Sigma\,\frac{r+1}{(z-n)^{2}} + \Sigma\,\frac{r+1}{(z-n_{1})^{2}}\right.
\]
\[
   \left.-\,\frac{p_{1}-1}{z-l}\,\frac{p_{1}-1}{z-l_{1}} - \cdots - \frac{q_{1}-1}{z-m}\,\frac{q_{1}-1}{z-m_{1}} - \cdots + \frac{r+1}{z-n}\,\frac{r+1}{z-n_{1}} + \cdots\right).
\]

\newpage
% ============================================================
% PDF p185 (printed 161)
% ============================================================
where it is to be observed that
\[
   \Sigma(p-1) + \Sigma(q-1) - \Sigma(r+1) = \delta - \sigma_{1} + \delta - \sigma_{2} - (\delta + \sigma_{3}) = \delta - \sigma_{1} - \sigma_{2} - \sigma_{3} = -2;
\]
consequently the function is diaphoric, and therefore curtate.

It is to be remarked that the function, although presenting itself in a form
unsymmetric in regard to the factors of $P$ and $Q$, and of $R$, is really symmetric
as regards the three sets of factors; this is obvious a priori, and it will be presently
verified.

\artno{23} For the calculation of the second term
\[
   x'^{2}\!\left(a, b, c \,.{}^{\!\cdot\!\cdot\!\cdot}\,\middle|\!\!\!\mid \frac{1}{x-a},\,\frac{1}{x-b},\,\frac{1}{x-c}\right)^{\!2}
\]
we have
\[
   f(x-a),\ g(x-b),\ h(x-c) = \Omega P,\ \Omega Q,\ \Omega R,
\]
where $\Omega$ is a determinate function of $z$; hence
\[
   \frac{x'}{x-a},\ \frac{x'}{x-b},\ \frac{x'}{x-c} = \frac{P'}{P} + \frac{\Omega'}{\Omega},\ \frac{Q'}{Q} + \frac{\Omega'}{\Omega},\ \frac{R'}{R} + \frac{\Omega'}{\Omega}.
\]
Then substituting these values, by reason that the function is diaphoric, the terms
in $\dfrac{\Omega'}{\Omega}$ disappear, and we have
\[
   x'^{2}\!\left(a, b, c \,.{}^{\!\cdot\!\cdot\!\cdot}\,\middle|\!\!\!\mid \frac{1}{x-a},\,\frac{1}{x-b},\,\frac{1}{x-c}\right)^{\!2}
   = \left(a, b, c \,.{}^{\!\cdot\!\cdot\!\cdot}\,\middle|\!\!\!\mid \frac{P'}{P},\,\frac{Q'}{Q},\,\frac{R'}{R}\right)^{\!2},
\]
which is
\[
   = \left(a, b, c \,.{}^{\!\cdot\!\cdot\!\cdot}\,\middle|\!\!\!\mid \Sigma\,\frac{p}{z-l},\,\Sigma\,\frac{q}{z-m},\,\Sigma\,\frac{r}{z-n}\right)^{\!2}.
\]

We have $\Sigma p = \Sigma q = \Sigma r$, $= \delta$: and hence by what precedes, this function, considered as
a function of the inverts $\dfrac{1}{z-l}$, \&c., is diaphoric, and therefore curtate.

\artno{24} We have therefore
\[
   \{x,\,z\} + x'^{2}\!\left(a, b, c \,.{}^{\!\cdot\!\cdot\!\cdot}\,\middle|\!\!\!\mid \frac{1}{x-a},\,\frac{1}{x-b},\,\frac{1}{x-c}\right)^{\!2}
\]
\[
   = \tfrac{1}{2}\!\left[\Sigma\,\frac{p-1}{(z-l)^{2}} + \Sigma\,\frac{q-1}{(z-m)^{2}} - \Sigma\,\frac{r+1}{(z-n)^{2}}\right]
\]
\[
   - \tfrac{1}{2}\!\left[\Sigma\,\frac{p-1}{z-l} + \Sigma\,\frac{q-1}{z-m} - \Sigma\,\frac{r+1}{z-n}\right]^{\!2}
   + \left(a, b, c \,.{}^{\!\cdot\!\cdot\!\cdot}\,\middle|\!\!\!\mid \Sigma\,\frac{p}{z-l},\,\Sigma\,\frac{q}{z-m},\,\Sigma\,\frac{r}{z-n}\right)^{\!2},
\]
where the whole function on the right-hand side is curtate.

\newpage
% ============================================================
% PDF p186 (printed 162)
% ============================================================
\artno{25} We have to bring the function on the right-hand side into the reduced form
\[
   \frac{a}{(z-a)^{2}} + \cdots + \frac{A}{z-a} + \cdots,
\]
for the purpose of getting rid of the non-essential inverts (if any).

We write
\[
   \Sigma\,\frac{p_{1}-1}{(z-l)^{2}} = \frac{p-1}{(z-l)^{2}} + \frac{p_{1}-1}{(z-l_{1})^{2}} + \cdots
\]
\[
   = \frac{p-1}{(z-l)^{2}} + \Sigma_{1}\,\frac{p_{1}-1}{(z-l_{1})^{2}};
\]
viz.\ $z-l$ here denotes any particular factor, and $z-l_{1}$ represents any other factor of
the same set; and so in other like cases.

\artno{26} The whole coefficient of $\dfrac{1}{(z-l)^{2}}$ is
\[
   = -(p-1) - \tfrac{1}{2}(p-1)^{2} + ap^{2},\ \ = \tfrac{1}{2}(1 - p^{2}) + ap^{2};
\]
an expression which, regarded as a function of $a$ and $p$, is represented by $(ap)$: the
parentheses are used only to avoid ambiguity, and are omitted when $p$ is a number,
thus $a1 = a$, $a2 = -\tfrac{3}{2} + 4a$, and so in other cases.

\artno{27} The whole term in $\dfrac{1}{z-l}$ comes from
\[
   -\frac{p-1}{z-l}\!\left[\Sigma_{1}\,\frac{p_{1}-1}{z-l_{1}} + \Sigma\,\frac{q-1}{z-m} - \Sigma\,\frac{r+1}{z-n}\right]
\]
\[
   + \frac{2ap}{z-l}\!\left[\Sigma_{1}\,\frac{p}{z-l_{1}} + \tfrac{1}{2}(-a-b+c)\Sigma\,\frac{q}{z-m} + \tfrac{1}{2}(-a+b-c)\Sigma\,\frac{r}{z-n}\right].
\]
Viz.\ each term such as $-\dfrac{p_{1}-1}{z-l_{1}}\cdot\dfrac{1}{z-l}$ is to be replaced by $\dfrac{1}{z-l\,.\,z-l_{1}}\!\left(\dfrac{1}{z-l} - \dfrac{1}{z-l_{1}}\right)$, giving
rise to the term $-\dfrac{p_{1}-1}{l-l_{1}}\,\dfrac{1}{z-l}$, or contributing the term $-\dfrac{1}{l-l_{1}}$ to the coefficient of $\dfrac{1}{z-l}$.
The whole coefficient thus is
\[
   = -(p-1)\!\left(\Sigma_{1}\,\frac{p_{1}-1}{l-l_{1}} + \Sigma\,\frac{q-1}{l-m} - \Sigma\,\frac{r+1}{l-n}\right)
\]
\[
   + 2ap\,\Sigma_{1}\,\frac{p}{l-l_{1}} + p(-a-b+c)\,\Sigma\,\frac{q}{l-m} + p(-a+b-c)\,\Sigma\,\frac{r}{l-n}.
\]

\artno{28} Suppose first that $z-l$ is a multiple factor of $P$, viz.\ a factor with an index $p$
greater than 1: then, for $z = l$, we have $Q + R = 0$, $Q' + R' = 0$, and thence $\dfrac{Q}{R} = \dfrac{Q'}{R'}$.
We have therefore $\Sigma\,\dfrac{q}{l-m} = \Sigma\,\dfrac{r}{l-n}$; that is,
\[
   p(-a-b+c)\,\Sigma\,\frac{q}{l-m} + p(-a+b-c)\,\Sigma\,\frac{r}{l-n}
\]

\newpage
% ============================================================
% PDF p187 (printed 163)
% ============================================================
moreover, in the top line, the terms $\Sigma\,\dfrac{q}{l-m}$ and $-\Sigma\,\dfrac{r}{l-n}$ destroy each other. The
whole coefficient of $\dfrac{1}{z-l}$, when $z-l$ is a multiple factor of $P$, thus is
\[
   = -(p-1)\!\left(\Sigma_{1}\,\frac{p_{1}-1}{l-l_{1}} + \Sigma\,\frac{1}{l-m} - \Sigma\,\frac{1}{l-n}\right)
\]
\[
   + 2ap\,\Sigma\,\frac{p_{1}}{l-l_{1}} - ap\,\Sigma\!\left(\frac{q}{l-m} + \Sigma\,\frac{r}{l-n}\right),
\]
a form which is now symmetrical in regard to the inverts $\dfrac{1}{l-m}$ and $\dfrac{1}{l-n}$.

\artno{29} The value just obtained must be equal to
\[
   (1 - p^{2} + 2ap)\!\left(\Sigma\,\frac{p_{1}-1}{l-l_{1}} + \tfrac{1}{2}\Sigma\,\frac{r-1}{l-n} - \Sigma\,\frac{1}{l-l_{1}}\right);
\]
viz.\ comparing the two forms and reducing, they will be identical if only
\[
   (1 - p + 2ap)\!\left\{\Sigma\,\frac{p+p_{1}}{l-l_{1}} - \Sigma\,\frac{\tfrac{1}{2}(1+p)q - p}{l-m} - \Sigma\,\frac{\tfrac{1}{2}(1+p)r - p}{l-n}\right\} = 0,
\]
and it can be shown that the function inside the [ ] is in fact $=0$.

\artno{30} We have, as before, $\Sigma\,\dfrac{q}{l-m} = \Sigma\,\dfrac{r}{l-n}$; or writing each of these quantities $= \Phi$,
the equation to be verified is
\[
   \Sigma\,\frac{p+p_{1}}{l-l_{1}} = (p+1)\,\Phi - p\,\Sigma\,\frac{1}{l-m} - p\,\Sigma\,\frac{1}{l-n}.
\]
We have
\[
   \frac{P'}{P} = \frac{P}{z-l} + \Sigma\,\frac{P_{1}}{z-l_{1}} = \frac{X_{1}}{X},
\]
that is,
\[
   \Sigma\,\frac{P_{1}}{l-l_{1}} = \left[\frac{X_{1}}{X} - \frac{p}{z-l}\right]\text{ for } z = l,
\]
\[
   = \left[\frac{X_{1}(z-l) - p X}{X(z-l)}\right].
\]
The first derived function of the numerator is $X_{1}'(z-l) + X_{1} - pX'$, which for
$z = l$ is $X_{1} - pX'$, which is $=0$; and, for the denominator, it is $X'(z-l) + X$, which
is also $=0$. Passing to the second derived functions, we find

\newpage
% ============================================================
% PDF p188 (printed 164)
% ============================================================
we find in like manner
\[
   \Sigma\,\frac{1}{l-l_{1}} = \tfrac{1}{2}\,\frac{X''}{X'},
\]
and we thence obtain ($z$ being always $=l$)
\[
   \Sigma\,\frac{p+p_{1}}{z-l_{1}} = \frac{X_{1}'}{X'},
\]
so that the equation to be verified becomes
\[
   \frac{X_{1}'}{X'} = (p+1)\,\Phi - p\,\Sigma\,\frac{1}{l-m} - p\,\Sigma\,\frac{1}{l-n}.
\]

\artno{31} But from the equation $\Theta = PQ' - P'Q$, $= KP_{0}Q_{0}R_{0}$, we find $XY_{1} - X_{1}Y = KR_{0}$,
and then, differentiating, $XY_{1}' + X'Y_{1} - X_{1}'Y - X_{1}Y' = KR_{0}'$: writing in these equations
$z = l$, they become
\[
   -X_{1}Y = KR_{0},
\]
\[
   X'Y_{1} - X_{1}'Y = KR_{0}',
\]
so that, dividing the second by the first,
\[
   -\frac{X'}{X_{1}}\,\frac{Y_{1}}{Y} + \frac{X_{1}'}{X_{1}} = \frac{R_{0}'}{R_{0}};
\]
or, recollecting that $X_{1} = pX'$ and $\dfrac{Y_{1}}{Y} = \dfrac{Q'}{Q}$, we have
\[
   \frac{X_{1}'}{X'} = p\!\left(\frac{R_{0}'}{R_{0}} - \frac{Y'}{Y}\right) + \frac{Q'}{Q},
\]
that is,
\[
   \frac{X_{1}'}{X'} = p\!\left(\Sigma\,\frac{r-1}{z-n} - \Sigma\,\frac{1}{z-m}\right) + \Sigma\,\frac{q}{l-m},
\]
\[
   = (p+1)\,\Phi - p\,\Sigma\,\frac{1}{l-m} - p\,\Sigma\,\frac{1}{l-n},
\]
the required relation.

\artno{32} The result is that, $z-l$ being a multiple factor of $P$, the coefficient of the
term $\dfrac{1}{z-l}$ is
\[
   = (1 - p^{2} + 2ap)\!\left\{\Sigma\,\frac{p_{1}-1}{l-l_{1}} + \tfrac{1}{2}\Sigma\,\frac{r-1}{l-n} - \Sigma\,\frac{1}{l-l_{1}}\right\}
\]
\[
   = 2\,(ap)\!\left[\tfrac{1}{2}\!\left(\frac{Q'}{Q} + \frac{R'}{R}\right) - \frac{Y'}{Y} - \frac{Z'}{Z} + \frac{X''}{X'}\right].
\]

\artno{33} In the case where $z-l$ is a simple factor of $P$ we have $p=1$, and the
coefficient is
\[
   = \Sigma\Sigma'\,\frac{p_{1}}{l-l_{1}} + (-a-b+c)\,\Sigma\,\frac{q}{l-m} + (-a+b-c)\,\Sigma\,\frac{r}{l-n}
\]
\[
   = a\!\left(2\Sigma'\,\frac{p_{1}}{l-l_{1}} - \Sigma\,\frac{q}{l-m} - \Sigma\,\frac{r}{l-n}\right) - (b-c)\!\left(\Sigma\,\frac{q}{l-m} - \Sigma\,\frac{r}{l-n}\right).
\]

\newpage
% ============================================================
% PDF p189 (printed 165)
% ============================================================
\artno{34} Of course the formulae for the coefficients of $\dfrac{1}{(z-l)^{2}}$ and $\dfrac{1}{z-l}$ give at once,
by a mere change of letters, those for the coefficients of $\dfrac{1}{(z-m)^{2}}$, $\dfrac{1}{z-m}$, and
$\dfrac{1}{(z-n)^{2}}$, $\dfrac{1}{z-n}$; and the function in question,
\[
   \{x,\,z\} + x'^{2}\!\left(a, b, c \,.{}^{\!\cdot\!\cdot\!\cdot}\,\middle|\!\!\!\mid \frac{1}{x-a},\,\frac{1}{x-b},\,\frac{1}{x-c}\right)^{\!2},
\]
is now obtained in the required form
\[
   = \frac{(ap)}{(z-l)^{2}} + \cdots + \frac{(bq)}{(z-m)^{2}} + \cdots + \frac{(cr)}{(z-n)^{2}} + \cdots + \frac{A}{z-l} + \cdots + \frac{B}{z-m} + \cdots + \frac{C}{z-n} + \cdots,
\]
where $(ap)$ denotes $\tfrac{1}{2}(1 - p^{2}) + ap^{2}$, and the like for $(bq)$ and $(cr)$; and where, $z-l$
being a multiple factor of $P$, the coefficient $A$ contains the factor $(ap)$, and similarly
for $B$ and $C$.

\artno{35} Suppose that the coefficients $a$, $b$, $c$ are no one of them $=0$; we have
$a1, =a$, which does not vanish; that is, $z-l$ being a simple factor of $P$, the
expression contains $\dfrac{1}{(z-l)^{2}}$, or the invert $\dfrac{1}{z-l}$ is essential: and similarly, $z-m$ being
a simple factor of $Q$, or $z-n$ a simple factor of $R$, the inverts $\dfrac{1}{z-m}$ and $\dfrac{1}{z-n}$
are essential. But for $z-l$ a multiple factor of $P$, the coefficient $(ap)$ of the term $\dfrac{1}{(z-l)^{2}}$
may vanish, viz.\ this will be the case if $a = \tfrac{1}{2}\!\left(1 - \dfrac{1}{p^{2}}\right)$; and, when this is so, the
coefficient $A$ of the corresponding term $\dfrac{1}{z-l}$ also vanishes; that is, $\dfrac{1}{z-l}$ is a non-essential invert. And similarly for any multiple factor $z-m$ of $Q$ or $z-n$ of $R$, the
invert $\dfrac{1}{z-m}$ or $\dfrac{1}{z-n}$ may be non-essential.

\artno{36} If $P$, $Q$, $R$ contain each of them only multiple factors of the same index,
say of the indices $p$, $q$, $r$ for the three functions respectively, viz.\ if the functions
are $P(\Pi(z-l))^{p}$, $G(\Pi(z-m))^{q}$, $H(\Pi(z-n))^{r}$, the result contains only the six terms
written for $a$, $b$, $c$ any given values of them $=\tfrac{1}{2}\!\left(1 - \dfrac{1}{p^{2}}\right),\ \tfrac{1}{2}\!\left(1 - \dfrac{1}{q^{2}}\right),\ \tfrac{1}{2}\!\left(1 - \dfrac{1}{r^{2}}\right)$ respectively
the result is $=0$: viz.\ we then have
\[
   \{x,\,z\} + x'^{2}\!\left(a, b, c \,.{}^{\!\cdot\!\cdot\!\cdot}\,\middle|\!\!\!\mid \frac{1}{x-a},\,\frac{1}{x-b},\,\frac{1}{x-c}\right)^{\!2} = 0,
\]
or we in fact have, for the values of $a$, $b$, $c$ in question,
\[
   f(x-a) : g(x-b) : h(x-c) = P : Q : R
\]
of this differential equation of the third order.

\newpage
% ============================================================
% PDF p190 (printed 166)
% ============================================================
\artno{37} The reasoning applies directly to lines 2, 3, 4, 5 of the PQR-Table; and
with a slight variation to line 1; viz.\ here the factors of $R$ ($=-z^{n}+1$) are all simple
factors, but in virtue of $c = 0$ and $n = b$, the corresponding inverts disappear, and, the
other inverts also disappearing, the value of the function in $\{x,\,z\}$ becomes $\equiv 0$.

Thus line 3 shows that the function $x$, determined by
\[
   f(x-a) : g(x-b) : h(x-c) = (z^{4} + 2\sqrt{-3}\,z^{2} + 1)^{3} : -12\sqrt{-3}\,(z^{4} - 1)^{2} : -(z^{4} - 2\sqrt{-3}\,z^{2} + 1)^{3},
\]
satisfies
\[
   \{x,\,z\} + x'^{2}\!\left(\tfrac{4}{9},\,\tfrac{3}{8},\,\tfrac{4}{9} \,.{}^{\!\cdot\!\cdot\!\cdot}\,\middle|\!\!\!\mid \frac{1}{x-a},\,\frac{1}{x-b},\,\frac{1}{x-c}\right)^{\!2} = 0,
\]
and so for any other of the five lines.

\artno{38} The indices of the factors of $P$, $Q$, $R$ may be such that, for proper values
of the coefficients $a$, $b$, $c$, there are in all only three essential inverts, say $\dfrac{1}{x-a_{1}}$,
$\dfrac{1}{x-b_{1}}$, $\dfrac{1}{x-c_{1}}$, belonging to the three functions $P$, $Q$, $R$ respectively, or it may be
two, or three, of them to the same function. When this is so, the function of those
inverts is, by what precedes, a curtate function, and it is consequently a function
\[
   \left(a_{1},\,b_{1},\,c_{1} \,.{}^{\!\cdot\!\cdot\!\cdot}\,\middle|\!\!\!\mid \frac{1}{x-a_{1}},\,\frac{1}{x-b_{1}},\,\frac{1}{x-c_{1}}\right)^{\!2},
\]
where $a_{1}$, $b_{1}$, $c_{1}$ are the values of the three which do not vanish in the series of
expressions $(ap)$, $(bq)$, $(cr)$.

The remaining lines (III, V, VII, VIII) and IX to XV of the PQR-Table give
such values of $P$, $Q$, $R$, the values of $(a, b, c)$ and the calculation of the values of
$(ap)$, $(bq)$, $(cr)$ is shown by the corresponding lines of the Annex. And we have thus
values of $x$ determined by the equations
\[
   f(x-a) : g(x-b) : h(x-c) = P : Q : R,
\]
and giving
\[
   \{x,\,z\} + x'^{2}\!\left(a, b, c \,.{}^{\!\cdot\!\cdot\!\cdot}\,\middle|\!\!\!\mid \frac{1}{x-a},\,\frac{1}{x-b},\,\frac{1}{x-c}\right)^{\!2} = \left(a_{1},\,b_{1},\,c_{1} \,.{}^{\!\cdot\!\cdot\!\cdot}\,\middle|\!\!\!\mid \frac{1}{x-a_{1}},\,\frac{1}{x-b_{1}},\,\frac{1}{x-c_{1}}\right)^{\!2}.
\]

\artno{39} For instance, from line IX we have
\[
   f(x-a) : g(x-b) : h(x-c) = (z-4)^{3} : -(z-1)(z+8)^{2} : 27 z^{2}\!\left(1 - \frac{z}{\infty}\right),
\]
the values of $(a, b, c)$ are $\tfrac{4}{9},\ \tfrac{3}{8},\ \tfrac{12}{25}$; and since $P$, $Q$, $R$ contain factors with the
exponents 3; 1, 2; and 1, 2 respectively, the coefficients which present themselves
on the right-hand side are
\[
   a3;\ \ b1,\ b2;\ \ c1,\ c2,
\]

\newpage
% ============================================================
% PDF p191 (printed 167)
% ============================================================
which are
\[
   = 0;\ \ \frac{3}{8},\ 0;\ \ \frac{12}{25},\ \frac{21}{50},\ \text{respectively.}
\]
Hence writing $a_{1}$, $b_{1}$, $c_{1} = \dfrac{3}{8},\ 0;\ \dfrac{12}{25},\ \dfrac{21}{50}$, the corresponding inverts are $\dfrac{1}{x-1},\ \dfrac{1}{x-\infty}$,
$\dfrac{1}{x}$; and the result is
\[
   \{x,\,z\} + x'^{2}\!\left(\frac{4}{9},\,\frac{3}{8},\,\frac{12}{25} \,.{}^{\!\cdot\!\cdot\!\cdot}\,\middle|\!\!\!\mid \frac{1}{x-a},\,\frac{1}{x-b},\,\frac{1}{x-c}\right)^{\!2}
   = \left(\frac{3}{8},\,\frac{12}{25},\,\frac{21}{50} \,.{}^{\!\cdot\!\cdot\!\cdot}\,\middle|\!\!\!\mid \frac{1}{x-1},\,\frac{1}{x-\infty},\,\frac{1}{x}\right)^{\!2}.
\]

\artno{40} It is hardly necessary to remark that an expression
\[
   \left(a_{1},\,b_{1},\,c_{1} \,.{}^{\!\cdot\!\cdot\!\cdot}\,\middle|\!\!\!\mid \frac{1}{x-a_{1}},\,\frac{1}{x-b_{1}},\,\frac{1}{x-c_{1}}\right)^{\!2}
\]
in fact denotes
\[
   \frac{a_{1}}{(x-a_{1})^{2}} + \frac{b_{1}}{(x-b_{1})^{2}} + \frac{-a_{1}-b_{1}+c_{1}}{(x-a_{1})(x-b_{1})} + \cdots
\]

The particular form of the $x$ inverts is immaterial; we could by a general linear
transformation upon the $x$ make them to be $\dfrac{1}{x-a_{1}},\ \dfrac{1}{x-b_{1}},\ \dfrac{1}{x-c_{1}}$ with the $(a_{1}, b_{1}, c_{1})$
arbitrary; or we can give to the $a_{1}, b_{1}, c_{1}$ any particular values we please: there
would be a propriety in making the inverts to be in every case (as in the foregoing
example) $\dfrac{1}{x},\ \dfrac{1}{x-\infty},\ \dfrac{1}{x-1}$; but the numerical work would be troublesome, and it is
not worth while to effect it.

\artno{41} The conclusion is that lines (III, V, VII, VIII) and IX to XV of the
PQR-Table, give, for determinate values of $(a, b, c)$ and $(a_{1}, b_{1}, c_{1})$, solutions
\[
   f(x-a) : g(x-b) : h(x-c) = P : Q : R
\]
of the equation
\[
   \{x,\,z\} + x'^{2}\!\left(a, b, c \,.{}^{\!\cdot\!\cdot\!\cdot}\,\middle|\!\!\!\mid \frac{1}{x-a},\,\frac{1}{x-b},\,\frac{1}{x-c}\right)^{\!2}
   = \left(a_{1}, b_{1}, c_{1} \,.{}^{\!\cdot\!\cdot\!\cdot}\,\middle|\!\!\!\mid \frac{1}{x-a_{1}},\,\frac{1}{x-b_{1}},\,\frac{1}{x-c_{1}}\right)^{\!2},
\]
where $a$, $b$, $c$, $a_{1}$, $b_{1}$, $c_{1}$ are or can be made arbitrary, but without any real gain of
generality herein. This is the Differential Equation $\{x,\,z\}$.

\artno{42} Recurring to the results from the Arabic lines of the PQR-Table, but for
convenience writing $s$ instead of $x$, we have
\[
   f(s-a) : g(s-b) : h(s-c) = P : Q : R,
\]
where $P$, $Q$, $R$ are now functions of $z$, a solution of
\[
   [s,\,z] + \left(\frac{ds}{dz}\right)^{\!2}\!\left(a, b, c \,.{}^{\!\cdot\!\cdot\!\cdot}\,\middle|\!\!\!\mid \frac{1}{s-a},\,\frac{1}{s-b},\,\frac{1}{s-c}\right)^{\!2} = 0.
\]
But we have
\[
   \{s,\,x\} = -\left(\frac{dx}{ds}\right)^{\!2}[s,\,z],
\]

\newpage
% ============================================================
% PDF p192 (printed 168)
% ============================================================
and the foregoing is therefore a solution of
\[
   [s,\,x] = \left(a, b, c \,.{}^{\!\cdot\!\cdot\!\cdot}\,\middle|\!\!\!\mid \frac{1}{s-a},\,\frac{1}{s-b},\,\frac{1}{s-c}\right)^{\!2},
\]
a differential equation of the third order. This is the Differential Equation $\{s,\,x\}$.

\artno{43} From the Roman lines, if we assume
\[
   f(s-a) : g(s-b) : h(s-c) = P : Q : R,
\]
where $\mathfrak{P}$, $\mathfrak{Q}$, $\mathfrak{R}$ are functions of $z$, not the same functions than $P$, $Q$, $R$ are of $s$,
since they belong to a different line of the Table: we have, as before,
\[
   [x,\,z] + \left(\frac{dx}{dz}\right)^{\!2}\!\left(a, b, c \,.{}^{\!\cdot\!\cdot\!\cdot}\,\middle|\!\!\!\mid \frac{1}{x-a},\,\frac{1}{x-b},\,\frac{1}{x-c}\right)^{\!2}
   = \left(a_{1}, b_{1}, c_{1} \,.{}^{\!\cdot\!\cdot\!\cdot}\,\middle|\!\!\!\mid \frac{1}{x-a_{1}},\,\frac{1}{x-b_{1}},\,\frac{1}{x-c_{1}}\right)^{\!2}.
\]

\artno{44} We may combine any such result with a properly selected result of the
preceding system, the two results being such that $(a, b, c)$ have the same values in
each of them. (See as to this the foot-note referring to the Annex to the PQR-Table.) The last equation then becomes
\[
   [x,\,z] + \left(\frac{dx}{dz}\right)^{\!2}[s,\,x] = \left(a_{1}, b_{1}, c_{1} \,.{}^{\!\cdot\!\cdot\!\cdot}\,\middle|\!\!\!\mid \frac{1}{x-a_{1}},\,\frac{1}{x-b_{1}},\,\frac{1}{x-c_{1}}\right)^{\!2},
\]
or since
\[
   [x,\,z] + \left(\frac{dx}{dz}\right)^{\!2}[s,\,x] = [s,\,z],
\]
this is
\[
   [s,\,z] = \left(a_{1}, b_{1}, c_{1} \,.{}^{\!\cdot\!\cdot\!\cdot}\,\middle|\!\!\!\mid \frac{1}{x-a_{1}},\,\frac{1}{x-b_{1}},\,\frac{1}{x-c_{1}}\right)^{\!2},
\]
the corresponding relation between $x$, $z$ being of course obtained by the elimination
of $s$ from the two sets of equations
\[
   f(s-a) : g(s-b) : h(s-c) = P : Q : R\ \text{ and }\ f(x-a) : g(x-b) : h(x-c) = \mathfrak{P} : \mathfrak{Q} : \mathfrak{R},
\]
viz.\ the required relation is
\[
   P : Q : R = \mathfrak{P} : \mathfrak{Q} : \mathfrak{R},
\]
where $P$, $Q$, $R$ are functions of $s$; $\mathfrak{P}$, $\mathfrak{Q}$, $\mathfrak{R}$ functions of $x$; and, in virtue of
\[
   P + Q + R = 0,\ \ \mathfrak{P} + \mathfrak{Q} + \mathfrak{R} = 0,
\]
the relations between equivalents to a single equation between $x$ and $s$. And writing
finally $x$ in place of $z$, that is, now considering $\mathfrak{P}$, $\mathfrak{Q}$, $\mathfrak{R}$ as functions of $x$, we have
\[
   \mathfrak{P} : \mathfrak{Q} : \mathfrak{R} = P : Q : R
\]
as a solution of
\[
   [s,\,x] = \left(a_{1}, b_{1}, c_{1} \,.{}^{\!\cdot\!\cdot\!\cdot}\,\middle|\!\!\!\mid \frac{1}{x-a_{1}},\,\frac{1}{x-b_{1}},\,\frac{1}{x-c_{1}}\right)^{\!2},
\]
a differential equation of the third order of the foregoing form, $[s,\,x] = $ given function
of $x$, but with different values of the coefficients $(a_{1}, b_{1}, c_{1})$ instead of $(a, b, c)$.

\newpage
% ============================================================
% PDF p193 (printed 169)
% ============================================================
\artno{45} It thus appears that there are in all 16 sets of values of $(a, b, c)$ for
which the equation is solved, viz.\ the 16 sets of values are shown in the right-hand
column of the Annex. For greater clearness I exhibit the integral equations
as follows:

\begin{center}
\begin{tabular}{|c|l|c|l|}
\hline
& Functions of $z$. & & Functions of $x$. \\
\hline
1 & $f(x-a) : g(x-b) : h(x-c) = P : Q : R\ (1)$ & $=$ & Polygon \\
I & ,, & $=$ & Double Pyramid \\[2pt]
II & $4z : -(z+1)^{2} : (z-1)^{2}\ (2)$ & $=$ & Tetrahedron \\
III & ,, & $=$ & ,, \\[2pt]
IV & $f(x-a) : g(x-b) : h(x-c) = (z-1)^{2} : (z+1)^{2}\ (4)$ & $=$ & Cube and Octahedron \\
V & ,, & $=$ & ,, \\[2pt]
VI & $f(x-a) : g(x-b) : h(x-c) = (4z + 1)^{2} : (z-1)^{2}\ (4)$ & $=$ & Dodecahedron and Icosahedron \\
VII & ,, & $=$ & ,, \\
VIII & $(z-1)^{2} : -(z+1)^{2} : 4z\ (2)$ & $=$ & ,, \\
IX & $P : Q : R\ (\text{IX})$ & $=$ & ,, \\
X & ,, & $=$ & ,, \\
XI & ,, & $=$ & ,, \\
XII & ,, & $=$ & ,, \\
XIII & ,, & $=$ & ,, \\
XIV & ,, & $=$ & ,, \\
XV & ,, & $=$ & ,, \\
\hline
\end{tabular}
\end{center}

The values of the $P$, $Q$, $R$ as functions of $z$, are taken out of the
$PQR$-Table: only in the lines III, V, VII, VIII, where $P$, $Q$, $R$ are given in the right-hand column of the Annex. For greater clearness I exhibit the integral equations
as follows:
\[
   = 4x,\ \ -(x+1)^{2},\ \ (x-1)^{2},
\]
and where, as regards V and VIII, line is a transposition of $P$ and $R$; 1 have
inverted the actual values of the $\mathfrak{P}$-functions. (See as to this the foot-note referring
to the Annex.)

\bigskip

\begin{center}
\textit{The Schwarzian Theory. Art.\ Nos.\ 46 to 62.}
\end{center}

\medskip

\artno{46} Considering the foregoing equation
\[
   [s,\,x] = \left(a_{1},\,b_{1},\,c_{1} \,.{}^{\!\cdot\!\cdot\!\cdot}\,\middle|\!\!\!\mid \frac{1}{x-a_{1}},\,\frac{1}{x-b_{1}},\,\frac{1}{x-c_{1}}\right)^{\!2}
\]
as a particular case of the equation $[s,\,x] = $ Rational function of $x$, $= R(x)$ suppose,
then we have in 1, I, II, IV, VI solutions of the form $x = $ Rational function of $s$.

\newpage
% ============================================================
% PDF p194 (printed 170)
% ============================================================
Consider, in general, a solution of this form, $x = F(s)$ a rational function of $s$: then
$s$ is an irrational function of $x$, and if $a$, $b$ are any values of $s$, then $[a_{1},\,b_{1}] = R(x)$,
$[a_{2},\,b_{2}] = R(x)$; that is, $[a_{2},\,b_{2}] = [a_{1},\,b_{1}]$, and therefore (ante, No.\ 7) $a_{2} = \dfrac{a a_{1} + b}{c a_{1} + d}$. And
then $x = F(a_{2}) = F\!\left(\dfrac{a a_{1} + b}{c a_{1} + d}\right)$, $= F(a_{1})$: viz.\ $F(s)$ is a rational function of $s$, transformable into itself by the several homographic transformations of a group of
such transformations: viz.\ taking $s$ a group; and conversely to each homographic transformation of $s$ we have a rational function of $s$, transformable into itself by the several homographic transformations of a group of
such transformations; viz.\ taking $s$ as the parameter of a point, the points which are
between any two roots whatever of the equation $x = F(s)$ there exists a homographic
relation of the form in question. Further, if $s_{1}$ is a root, the others are
$\dfrac{a s_{1} + b}{c s_{1} + d}$ for the several homographic transformations $\dfrac{a s + b}{c s + d}$ of the group; and
transformable into itself by the several homographic transformations of a group of
such transformations occur in the form $x = F(s)$; conversely to each homographic transformation $\dfrac{a s + b}{c s + d}$ of the group correspond to the six elements of the
rational function $x = F(s)$ a Group; or, what is the same thing, to each homographic transformation $\dfrac{a s + b}{c s + d}$ of the group correspond to a Group; or, what is the same thing, to each homographic transformation $\dfrac{a s + b}{c s + d}$ of the group correspond to a Group; the rational function $x = F(s)$ is the function corresponding to the Group.

\artno{47} We may, in any equation between $x$ and $s$, consider $s$ as imaginary
variables $p + qi$ and $u + vi$ respectively; considering then $(p, q)$ and $(u, v)$ as rect-
angular coordinates of points in different planes, we have a first plane the locus of
the points $s$, and a second plane the locus of the points $x$: there is between the
two planes a correspondence which is in fact the correspondence of conformable
figures; to the infinitesimal element $ds$ drawn from a point $s$ of the first figure
corresponds an infinitesimal element $dx$ drawn from the corresponding point $x$ of the
second figure, these elements being in general connected by an equation of the
form $dx = \alpha\,ds$ involving $s$ but not $ds$; or, what is the same thing, the lengths of the
ratio of the lengths of the elements being the moduli of $\alpha$ ($a + bi$), and
the inclinations of the directions being the angle of $\alpha$ ($a + bi$). Passing to the
second derived element $d_{2}s$ of the same thing along $ds$, $d_{2}x$ the corresponding consecutive elements of the path of the point $s$, then the
ratio of the lengths of the elements $d_{2}x$, $d_{2}s$ are the corresponding consecutive elements of the path of the point $x$, the
ratio $d_{2}x : d_{2}s$ the corresponding consecutive elements of the path of the path of the point $s$, then the
ratio of the lengths of the elements $d_{2}s$ and $d_{2}x$ are the corresponding consecutive elements of the path of the point $x$. In particular, the consecutive elements $ds$
having all the same direction (that is, the same module $\lambda$), as $\lambda$ is a real
positive, has $s = a + bi$, $b$, $c$ are critical points of the kind in question. In fact, writing
$x - a = h$, the equation is of the form
\[
   x - a = i F'\,\exp i\,(\theta + \theta_{1}).
\]

\newpage
% ============================================================
% PDF p195 (printed 171)
% ============================================================
\artno{48} I consider the foregoing equation $[s,\,x] = R(x)$, where $R(x)$ is a rational
function, and is now taken to be a real function of $x$: we may assume $s' = i\rho'\,e^{i\theta}$,
where the accents denote differentiation in regard to $x$, and where $\rho'$, $\theta'$, and therefore also $\rho$, are real functions of $x$. We have
\[
   \frac{s''}{s'} = \frac{\rho''}{\rho'} + i\theta'',
\]
and thence
\[
   \frac{s'''}{s'} - \left(\frac{s''}{s'}\right)^{\!2} = \frac{\rho'''}{\rho'} - \left(\frac{\rho''}{\rho'}\right)^{\!2} + \frac{\theta'''}{\theta'} - \left(\frac{\theta''}{\theta'}\right)^{\!2} + i\theta''
\]
\[
   -\tfrac{1}{2}\!\left(\frac{s''}{s'}\right)^{\!2} = -\tfrac{1}{2}\!\left(\frac{\rho''}{\rho'}\right)^{\!2} - \tfrac{1}{2}\!\left(\frac{\theta''}{\theta'}\right)^{\!2} + \tfrac{1}{2}\theta'^{2} - i\frac{\rho''}{\rho'}\,i\theta'' - i\frac{\rho''\theta''}{\rho'},
\]
and thence
\[
   [s,\,x] = [\rho,\,x] + [\theta,\,x] + \tfrac{1}{2}\theta'^{2} - i\frac{\rho''\theta''}{\rho'\theta'} - i\frac{\rho''\theta''}{\rho'}.
\]

Putting this $= R(x)$, and assuming that $x$ is real, we have
\[
   [\rho,\,x] + [\theta,\,x] + \tfrac{1}{2}\theta'^{2} = R(x),\quad 0 = i\frac{\rho''\theta''}{\rho'\theta'}.
\]
The last equation gives $\rho''\theta'' = 0$, that is, $\theta'' = 0$, which gives $s' = 0$, and may be
disregarded; or else $\rho'' = 0$, therefore $\rho'$, a real constant, $= \gamma$ suppose, and $[\rho,\,x] = 0$;
hence for the solution of the equation $[s,\,x] = R(x)$, we have $s' = i\gamma'\,e^{i\theta}$, $\theta$ a real
quantity determined by $[\theta,\,x] + \tfrac{1}{2}\theta'^{2} = R(x)$; and then, integrating the equation for $s'$,
we have $s = a + \beta s + \gamma e^{i\theta}$, $a$, $\beta$, $\gamma$ real constants.

\artno{49} The conclusion is that, if $[s,\,x] = R(x)$, a real function of $x$, and if $x$ be
real, that is, if the point $x$ moves along a right line (say the $x$-line), then $s = a + \beta s + \gamma e^{i\theta}$
($\theta$, and the constants $a$, $\beta$, $\gamma$ being real), that is, the point $s$ moves in a circle,
coordinates of the center: $a$, $\beta$, and radius $= \gamma$.

\begin{center}
\textit{[Figure: Three circles labelled $A$, $B$, $C$ along the $x$-axis $\rho\!-\!\phi\!-\!q$, illustrating the conformal mapping of the real $x$-line to a $s$-plane consisting of arcs forming a triangular region.]}
\end{center}

\artno{50} Suppose $a$, $b$, $c$ are any real values of $x$ representing points $a$, $b$, $c$ on the
$x$-line; and $A$, $B$, $C$ any given imaginary values of $s$ representing points $A$, $B$, $C$

\newpage
% ============================================================
% PDF p196 (printed 172)
% ============================================================
in the $s$-plane: since $[s,\,x] = R(x)$ is a differential equation of the third order, the
integral contains three arbitrary constants, and we may imagine these so determined
that to the values $s = a, b, c$ shall correspond the values $s = A, B, C$ respectively.

If there is not on the $x$-line any critical point, as the point $x$ moves continuously along this line from the point $a$ towards the point $b$, the corresponding circle (in-
asmuch as $a$, $b$, $c$ and $A$, $B$, $C$ are corresponding points) must be the circle through
the three points $A$, $B$, $C$.

\artno{51} If however the points $a$, $b$, $c$ are critical points, such that the element $ds$
at the corresponding points $A$, $B$, $C$ are equal to multiples of $(dx)^{\lambda}$, $(dx)^{\mu}$, $(dx)^{\nu}$ re-
spectively, then on the $x$-line the point $a$ in the direction $ax$ to a passes continu-
ously along the line: the point $s$ at the point $A$ moves continuously along a circle, (in-
asmuch as $a$, $b$, $c$ and $A$, $B$, $C$ are corresponding points) must be the circle through
the three points $A$, $B$, $C$.

\artno{51} If however the points $a$, $b$, $c$ are critical points, such that the elements $ds$
only along the $s$-line: the point $a$, $b$, $c$ are critical points, such that the elements $ds$
along the $a$-line, $\mu\pi$, $\nu\pi$, of $(dx)^{\lambda}$, $(dx)^{\mu}$, $(dx)^{\nu}$ respectively. The arc $ABC$ becomes continuous
ously at $a$, $b$, $c$ but the point $s$ at the points $A$, $B$, $C$ must change abruptly through
the angles $\lambda\pi$, $\mu\pi$, $\nu\pi$, being real), that is, the point $s$ moves in a circle.

\artno{51} If however there is on the $x$-line a critical point, such that the element $ds$
only along the $x$-line. Pasing the values of $\lambda$, $\mu$, $\nu$, being real, and equal to the curvilinear triangle
that to the values $x = a, b, c$ shall correspond the values $s = A, B, C$. (In-
asmuch as $a$, $b$, $c$ and $A$, $B$, $C$ are corresponding points) must be the circle through
the three points $A$, $B$, $C$.

\artno{52} The foregoing equation
\[
   [s,\,x] = \left(a, b, c \,.{}^{\!\cdot\!\cdot\!\cdot}\,\middle|\!\!\!\mid \frac{1}{x-a},\,\frac{1}{x-b},\,\frac{1}{x-c}\right)^{\!2}
\]
where $a$, $b$ are real values of $x$ representing points $a$, $b$, $c$ on the
\[
   \frac{d}{dx}\log\frac{ds}{ds} = \frac{1 - \lambda}{h} + a_{1} + a_{2}h + a_{3}h^{2} + \cdots,
\]
which is satisfied. And similarly the other conditions are satisfied.
\[
   s = h k^{\lambda-1}(1 + b_{1}h + b_{2}h^{2} + \cdots), = k\rho\ \text{for shortness.}
\]
This is a particular integral, but we have from it the general integral
\[
   s = \frac{\alpha + \beta k\rho}{\gamma + \delta k\rho}.
\]

\noindent\small * Since there is no critical point on the $x$-line there can be no abrupt change of direction in the path
of $s$, that is, the path of $s$ cannot consist of circular arcs meeting at an angle: but it is in the next
further assumed that the path of $s$ cannot consist of different arcs of circles, the one continuing the other
without any abrupt change of direction.\normalsize

\newpage
% ============================================================
% PDF p197 (printed 173)
% ============================================================
If $A$ be the value of $s$ corresponding to a vanishing value of $h$, then $B = aA$, and we find
\[
   s = \frac{a + A\delta h\rho}{\gamma + \delta h\rho},\quad = \left(A + \frac{a}{\delta h\rho}\right)\!\left(1 + \frac{\gamma}{\delta h\rho}\right)^{\!-1},\quad = A + \frac{a - \gamma A}{\delta}\,\frac{1}{h\rho} + \cdots;
\]
viz.\ reducing $\dfrac{1}{\rho}$ to its principal term $h^{\lambda}$, and then writing $ds$, $ds$ for $s - A$, and $h\ (=x-a)$
respectively, we have $ds = K\,(dx)^{\lambda}$, or $s - a = \alpha$ is a critical point with the exponent $\lambda$;
and similarly $x = b$ and $x = c$ are critical points with the exponents $\mu$ and $\nu$ respectively.

\artno{53} Hence in the equation
\[
   [s,\,x] = \left(a, b, c \,.{}^{\!\cdot\!\cdot\!\cdot}\,\middle|\!\!\!\mid \frac{1}{x-a},\,\frac{1}{x-b},\,\frac{1}{x-c}\right)^{\!2},
\]
as the point $x$, passing successively through $a$, $b$, $c$, describes the $x$-line, the point $s$,
passing successively through $A$, $B$, $C$, describes the circular arcs meeting at the proper
triangle $ABC$. To points $s$ indefinitely near the $x$-line correspond points $x$ indefinitely
near to the boundary $ABC$ of the triangle $ABC$ on the $s$-plane. Similar reasoning shows
that the $s$-plane, supposing the boundary just spoken of, agrees fully with the $s$-plane in question, viz.\ the points $A$, $B$, $C$ are critical points of $x$, with the exponents $\lambda$, $\mu$, $\nu$,
where $a$, $b$, $c$, $\lambda$, $\mu$, $\nu$ are connected by the relation
\[
   f(x-a) : g(x-b) : h(x-c) = P : Q : R
\]
of this differential equation of the third order.

\artno{54} It is to be remarked, in regard to the values of the values of the values of the values of the values; and conversely, there are two values of $x$, viz.\
one in the upper half and the other in the lower half of the $x$-line, which correspond to the same
each of the points $s$ within the curvilinear triangle $ABC$. The two half-planes (above
and below) form, between them, two sets of triangles $ABC$, $\&c.,$ and $\&c.,$ respectively,
each of them which are equivalent functions of $x$. And we have thus the values $x$,
the points $A$, $B$, $C$ as the values, on opposite sides of the $s$-line, of the curvilinear triangles in the upper half-plane corresponds on opposite sides
of the $x$-line, correspond to two sets of triangles $ABC$, $\&c.,$ and $\&c.,$ respectively,
certain determinate values; and, in fact, dividing the surface of a sphere into
triangles $ABC$, $ABC$, $\&c.,$ alternately of the two kinds: corresponding to all these triangles in
the $s$-plane and the corresponding triangles in the $s$-plane are these
triangles; and hence the shaded half-plane corresponds to one set of them and
the unshaded triangles correspond to the other. To the values $x$, in the upper half of the $x$-line we have, $A$, $B$, $C$ corresponding to the foot-note for the values of $\frac{4}{5}$, $\frac{15}{50}$, etc., and there will be six considerations as these
last followers, in the Annex of the PQR-Table; and then showed a $a$ priori that the equation
\[
   [s,\,x] = \left(a, b, c \,.{}^{\!\cdot\!\cdot\!\cdot}\,\middle|\!\!\!\mid \frac{1}{x-a},\,\frac{1}{x-b},\,\frac{1}{x-c}\right)^{\!2}
\]
is algebraically integrable for these values of $a$, $b$, $c$; and only for these values, or
for values reducible to them.

\newpage
% ============================================================
% PDF p198 (printed 174)
% ============================================================
\artno{55} As an instance, take the double pyramid form: the integral equation is
\[
   f(x-a) : g(x-b) : h(x-c) = 4 x^{n} : -(x^{n} - 1)^{2} : (x^{n} + 1)^{2},
\]
or say
\[
   \frac{(x-a)(x-b)}{(x-b)(x-c)} = -\frac{(x^{n} - 1)^{2}}{(x^{n} + 1)^{2}};
\]
or if, for greater simplicity, we assume $a$, $b$, $c$, $= 1$, $0$, $\infty$, this is $x = \dfrac{(x^{n} - 1)^{2}}{(x^{n} + 1)^{2}}$, or say
$-(x^{n} - 1) = \sqrt{x}\,(x^{n} + 1)$, that is, $x^{n} = \dfrac{1 \mp \sqrt{x}}{1 \pm \sqrt{x}}$, a solution of the differential equation
\[
   [s,\,x] = \left(\frac{3}{8},\,\tfrac{1}{2}(1 - n^{-2}),\,\tfrac{3}{8} \,.{}^{\!\cdot\!\cdot\!\cdot}\,\middle|\!\!\!\mid \frac{1}{x - 1},\,\frac{1}{x - \infty},\,\frac{1}{x}\right)^{\!2}.
\]
In particular, if $n = 3$, we have $x = \dfrac{(s^{3} - 1)^{2}}{(s^{3} + 1)^{2}}$ or $s^{3} = \dfrac{1 \mp \sqrt{x}}{1 \pm \sqrt{x}}$, a solution of
\[
   [s,\,x] = \left(\frac{3}{8},\,\frac{4}{9},\,\frac{3}{8} \,.{}^{\!\cdot\!\cdot\!\cdot}\,\middle|\!\!\!\mid \frac{1}{x - 1},\,\frac{1}{x - \infty},\,\frac{1}{x}\right)^{\!2}.
\]

\artno{57} We have here the spherical surface divided by the equator and three meridians
into twelve triangles, each with the angles $\tfrac{1}{2}\pi,\,\tfrac{1}{3}\pi,\,\tfrac{1}{2}\pi$; and then, projecting from the
South pole on the plane of the equator, we have the annexed figure of the $s$-plane,

\begin{center}
\textit{[Figure: A central circle on the $s$-plane is divided by three radial lines meeting at the centre into six sectors; the perimeter is split by these and three meridional arcs into twelve curvilinear triangles, alternately shaded and unshaded. Labels $A, B, C, D$ mark four boundary points; another label sits at the centre and at $\infty$ outside the disk.]}
\end{center}

\noindent divided into 12 curvilinear triangles, each with these same angles $30^{\circ}, 90^{\circ}, 60^{\circ}$; the
plane is divided by the shading into two systems, each of 6 triangles. The figure
of the $s$-plane is by the $s$-line divided into two half-planes, one shaded, the other
unshaded; and we have on the line the point $c$ at $\infty$, $a$ at the origin, and $b$ at
the distance unity.

\newpage
% ============================================================
% PDF p199 (printed 175)
% ============================================================
\artno{58} Take $x$ real; then, if $x$ is positive and less than 1, $s^{3}$ is real and positive,
and we have for $s$ the infinitum half-lines at the inclinations $0^{\circ},\,180^{\circ},\,300^{\circ}$. If $x$ is real and negative,
or in particular greater than 1, viz.\ here the factors of $R\ (=-x^{n}+1)$ are all simple
factors, but in virtue of $c = 0$ and $n = b$, the corresponding inverts disappear, and, the
the form $\dfrac{1 - hi}{1 + hi}$, $= \cos\theta + i\sin\theta$; whence as $x$ is of the same form, or the locus of the
point $s$ is a circle radius unity. Writing $s^{3} = \dfrac{1 - \sqrt{x}}{1 + \sqrt{x}}$, and supposing that the point $x$
moves along the $x$-line through 1 through $\infty$ to $\infty$, and thence from $\infty$ through 0 to $b$,
the point $s$ describes the sides $BA,\, AC,\, CB$ of the shaded triangle marked $K$.

\artno{59} Suppose that the point $x$ is at $k$, in the shaded half-plane and at indefinitely
small distance $\epsilon$ from the $x$-line; in the shaded line in question, in the shaded triangle
$s\,(-i)$, we have $s^{3} = \dfrac{1 - \sqrt{x}}{1 + \sqrt{x}}$, and then $x = -\frac{1}{4}\epsilon(1 - i)$ nearly, and hence a value of $s$ is
$s = \dfrac{1 - ki}{1 + ki}$, $= \cos\theta + i\sin\theta$; whence $s^{3}$ is of the same form, or the locus of the
point $s$ is a circle radius unity. Writing $s^{3} = \dfrac{1 - \sqrt{x}}{1 + \sqrt{x}}$, and supposing that the point $x$
moves along the $x$-line $a$ through $b$ through $a$ to $\infty$, and to $\infty$ from $\infty$ to $b$, the shaded value $\sqrt{x}$ we have $s^{3} = \dfrac{1 - \epsilon(1 - i)}{1 + \epsilon(1 - i)}, = 1 - 2\epsilon(1 - i)$ nearly, and hence a value of $s$ is
$\epsilon = 1 - \tfrac{2}{3}\epsilon + \tfrac{2}{3}\epsilon i$ nearly, and we have to find from $a$ at $c$ in $K$, and $b$ at the situate on
the circle in question.

\artno{59} Suppose the point $x$ is at $k$, in the shaded half-plane, in the shaded triangle
small distance $\epsilon$ from the $x$-line, in the shaded line in question, then for the value
$x\ (-a) = h$, the shaded line $BA,\, AC,\, CB$ of the equation, which $\dfrac{1}{x-a}$ value
$\epsilon\,(1 - i)$, we have $s^{n} = \dfrac{1 - \epsilon(1 - i)}{1 + \epsilon(1 - i)}$, $= 1 - 2\epsilon(1 - i)$ nearly, and hence a value of $s$ is
$s = 1 - \tfrac{2}{3}\epsilon + \tfrac{2}{3}\epsilon i$, which belongs to the shaded triangle $K$; and similarly the other
triangles; and hence the shaded half-plane corresponds to the shaded triangles, and
the unshaded half-plane corresponds to the other.

\artno{60} Suppose the point $x$, instead of being in question of $a$, $b$, $c$, is in
the circle in question.

We have here $x - a = a\,(e^{i\theta_{1}} - e^{-i\theta_{1}}) = -2i a\,e^{i\frac{1}{2}\pi}\sin\tfrac{1}{2}\theta_{1};\,e^{i\theta_{1}} \cdot e^{i\theta_{2}}$, say
$x - a = a\,e^{i\theta_{1}} \cdot e^{i\theta_{2}}\,e^{i\frac{1}{2}\pi - \frac{1}{2}\theta_{2}}\,$ similarly
\[
   x - a = i P_{1}\,\exp\tfrac{1}{2}(\theta + \theta_{1}),
\]
Similarly $x - b = i Q_{1}\,\exp\tfrac{1}{2}(\theta + \theta_{2})$, and $x - c = i R_{1}\,\exp\tfrac{1}{2}(\theta + \theta_{3})$, where $P,\, Q,\, R$ denote $\sin\tfrac{1}{2}(\theta - \theta_{1})$, $\sin\tfrac{1}{2}(\theta - \theta_{2})$, $\sin\tfrac{1}{2}(\theta - \theta_{3})$ respectively. In like manner, $a - b$, $b - c$, $c - a$, $iH\exp\tfrac{1}{2}(\theta_{1} + \theta_{2})$, $iH\exp\tfrac{1}{2}(\theta_{2} + \theta_{3})$, $iH\exp\tfrac{1}{2}(\theta_{3} + \theta_{1})$, where $F,\, G,\, H$ denote $\sin\tfrac{1}{2}(\theta_{1} - \theta_{2})$, $\sin\tfrac{1}{2}(\theta_{2} - \theta_{3})$, $\sin\tfrac{1}{2}(\theta_{3} - \theta_{1})$ respectively.

\newpage
% ============================================================
% PDF p200 (printed 176)
% ============================================================
We have
\[
   -\frac{b - x \,.\, x - a \,.\, a - b}{x - a \,.\, x - b} = \frac{-FGH}{PQR}\,\exp i\tfrac{1}{2}\,(\theta_{1} + \theta_{2} + \theta_{3} - 3\theta),
\]
\[
   \frac{1}{b - c \,.\, x - a} = \frac{-1}{c^{2}IP}\,\exp i - \tfrac{1}{2}\,(\theta_{2} + \theta_{3} + \theta_{1} + \theta),
\]
with the like values for $\dfrac{1}{c - a \,.\, x - b}$ and $\dfrac{1}{a - b \,.\, x - c}$. Hence the right-hand side of
the equation is
\[
   = \frac{FGH}{PQR}\!\left(\frac{a}{PP} + \frac{b}{QQ} + \frac{c}{RR}\right)\exp i\,(-2\theta).
\]

\artno{62} Considering now the left-hand side of the equation, we have
\[
   [s,\,x] = \frac{1}{\!\left(\dfrac{dx}{d\theta}\right)^{\!2}}\,(\{s,\,\theta\} - \{x,\,\theta\});
\]
substituting for $x$ its value $= a + \beta i\,e^{i\theta} + \gamma e^{i\theta}$, this becomes
\[
   [s,\,x] = -\frac{1}{\gamma^{2}}\,e^{-i\theta}\,(\{s,\,\theta\} - \tfrac{1}{2}),
\]
that is,
\[
   = -\frac{1}{\gamma^{2}}\,(\{s,\,\theta\} - \tfrac{1}{2})\exp i\,(-2\theta).
\]

Assume $s = L + Mi + Ne^{i\theta}$, $L$, $M$ and $N$ constants; then using the accent to denote
differentiation in regard to $\theta$, we find without difficulty $[s,\,\theta] = (\Theta_{1} + \Theta_{2} + \tfrac{1}{2}\Theta^{2})$, and the
value of $[s,\,x]$ becomes
\[
   = -\frac{1}{\gamma^{2}}\!\left((\Theta_{1} + \Theta_{2} + \tfrac{1}{2}\Theta^{2}) - \tfrac{1}{2}\right)\exp i\,(-2\theta).
\]
Hence, substituting the values of the two sides of the equation, the imaginary
factor $\exp i\,(-2\theta)$ divides out, and the equation becomes
\[
   (\Theta_{1} + \Theta_{2} + \tfrac{1}{2}\Theta^{2}) = \frac{FGH}{PQR}\!\left(\frac{a}{PP} + \frac{b}{QQ} + \frac{c}{RR}\right),
\]
an equation, in which everything is real and which thus determines $\Theta$ as a real
function of $\theta$; and we have therefore the theorem in question.

\bigskip

\begin{center}
\textit{Connexion with the differential equation for the hypergeometric series. Art.\ Nos.\ 63 to 68.}
\end{center}

\medskip

\artno{63} Take $p$, $q$ given functions of $x$, and $y$ a function of $x$ determined by the
equation
\[
   \frac{d^{2}y}{dx^{2}} + p\,\frac{dy}{dx} + qy = 0;
\]

\end{document}
